\documentclass{aa}

\usepackage{txfonts}
\usepackage{orcidlink}

% suppress A&A journal header for arXiv preprint
\makeatletter
\renewcommand*\aa@journalname{arXiv-only}
\renewcommand*\aa@manuscriptname{}
\renewcommand*\aa@publishlink{}
\renewcommand*\aa@textidlineempty{} % TODO: set to arXiv:XXXX.XXXXX after submission
\makeatother

\newcommand{\crires}{CRIRES$^+$}
\newcommand{\um}{$\mu$m}
\newcommand{\kms}{km\,s$^{-1}$}
\newcommand{\vipere}{\texttt{vip\`ere}}
\newcommand{\esorex}{\texttt{esorex}}
\newcommand{\crtwores}{\texttt{cr2res}}

\begin{document}

\title{CRIRES$^+$ L- and M-band spectral archive with telluric-calibrated slit
 tilt and wavelength solutions}
\titlerunning{CRIRES$^+$ L/M-band spectral archive}

\author{Thomas Marquart\inst{1}\,\orcidlink{0000-0002-2244-9920}}
\authorrunning{T.\ Marquart}

\institute{Department of Physics and Astronomy, Uppsala University,
  Box 516, SE-75120 Uppsala, Sweden\\
  \email{thomas.marquart@astro.uu.se}}

\date{March 2026}

% TODOs
%
% * why only flats, no bpm, dark, etc
% * xcorr tried but no avail
% * check bet_Pic_M4318_2023-11-01_05255_2 after rerun
% * choice to neglect in detector-order

\abstract{
We present a uniformly reprocessed archive of all public \crires{} L-band
(2.8--4.2\,\um) and M-band (4.2--5.5\,\um) science observations obtained
between September 2021 and March 2025, totalling 11\,131 raw frames
across 16 wavelength settings. The reprocessing improves upon the standard
ESO pipeline by incorporating a slit tilt calibration derived from
telluric standard observations, per-observation flexure correction of the
spectral traces, and telluric correction using the \vipere{} forward
modelling tool. The resulting 5649 extracted AB nod-pair spectra, spanning
156 unique targets from 68 ESO programmes, are served through an
interactive web archive.
}

\keywords{infrared: general -- techniques: spectroscopic -- atmospheric effects --
methods: data analysis -- instrumentation: spectrographs}

\maketitle

%--------------------------------------------------------------------
\section{Introduction}
\label{sec:intro}

\crires{} is the upgraded cryogenic high-resolution infrared echelle
spectrograph at ESO's Very Large Telescope \citep{Dorn2023}. It provides
cross-dispersed spectroscopy from 0.95 to 5.3\,\um{} with a spectral
resolving power of $R \approx 100\,000$, recording cross-dispersed long-slit
spectra on 
three HAWAII2RG detectors
of $2048 \times 2048$ pixels each with 5--7 echelle orders per chip.

While the near-infrared ($YJHK$) modes of \crires{} are well supported by
the standard \crtwores{} pipeline \citep{cr2res} and included in ESO's
Phase~3 data releases, the thermal infrared L-band (2.8--4.2\,\um) and
M-band (4.2--5.5\,\um) settings are explicitly excluded from the Phase~3
archive. Two instrument-specific issues affect L/M reductions that the
standard pipeline does not address: (1) the slit tilt, which is
wavelength-dependent and significant in the thermal infrared, is not
calibrated for these bands; and (2) mechanical flexure of the instrument
shifts the spectral traces on the detector between calibration and science
observations, degrading extraction quality when uncorrected.

This paper describes a bulk reprocessing of all public \crires{} L/M
science data that addresses both issues and additionally provides telluric
correction. The reduced spectra are served through an interactive web
archive at \url{https://neon.physics.uu.se/crires-lm/}. The reduction
code is publicly available at
\url{https://github.com/ivh/CRIRES-LM}.


%--------------------------------------------------------------------
\section{Data}
\label{sec:data}

We queried the ESO archive for all \crires{} science frames with
wavelength settings in the L- and M-bands, covering the period from
September 2021 (first \crires{} science operations) through January 2025.
This yielded 11\,131 raw frames from 68 programmes observing 156 distinct
targets across 16 wavelength settings (Table~\ref{tab:settings}).

The observations use the standard \texttt{CRIRES\_spec\_obs\_AutoNodOnSlit}
template with ABBA nodding along the slit. Frames were paired into AB
nod pairs by greedy matching within each observing template sequence,
producing 5649 pairs across 412 combined observing sequences. An
additional 233 flat-field calibration epochs (deep flats with
NDIT\,$\geq$\,10) were retrieved for the same period.

\begin{table}
\caption{Wavelength settings and data volume.}
\label{tab:settings}
\centering
\begin{tabular}{lrrr}
\hline\hline
Setting & Range (nm) & Frames & Pairs \\
\hline
L3244 & 2869--3982 & 19 & 16 \\
L3262 & 2886--4003 & 972 & 498 \\
L3302 & 2923--4050 & 921 & 446 \\
L3340 & 2960--4082 & 765 & 399 \\
L3377 & 2843--4136 & 726 & 395 \\
L3412 & 2873--4177 & 61 & 36 \\
L3426 & 2886--4193 & 222 & 101 \\
\hline
M4187 & 3584--5504 & 76 & 40 \\
M4211 & 3381--5535 & 1118 & 578 \\
M4266 & 3427--5604 & 137 & 68 \\
M4318 & 3470--5154 & 1241 & 638 \\
M4368 & 3513--5211 & 4761 & 2398 \\
M4416 & 3552--5264 & 18 & 12 \\
M4461 & 3614--5314 & 74 & 39 \\
M4504 & 3870--5364 & 12 & 8 \\
M4519 & 3435--5370 & 8 & 4 \\
\hline
Total & & 11\,131 & 5649 \\
\hline
\end{tabular}
\end{table}


%--------------------------------------------------------------------
\section{Method}
\label{sec:method}

The reduction proceeds in four stages: slit tilt calibration from
telluric standard stars (\S\,\ref{sec:tilt}), wavelength calibration
update (\S\,\ref{sec:wavecal}), per-observation trace adjustment and
extraction (\S\,\ref{sec:extraction}), and telluric correction
(\S\,\ref{sec:tellcorr}).


%--------------------------------------------------------------------
\subsection{Slit tilt calibration}
\label{sec:tilt}

The \crires{} slit is not perfectly aligned with the detector columns;
the resulting tilt is wavelength-dependent and, if uncorrected, causes
spectral features to shift in wavelength as a function of position along
the slit. For nodding observations, where the spectrum is extracted at
two different slit positions, this introduces a systematic wavelength
offset between the A and B nod spectra. In the L and M bands, the slit
tilt ranges from about $+0.01$ to $-0.09$ pixels per spatial pixel,
corresponding to wavelength offsets of several \kms{} over the nod throw.

We calibrate the slit tilt using a dedicated telluric standard star for
each of the 16 wavelength settings (Table~\ref{tab:tellstds}). These are
hot B- and O-type stars whose intrinsic spectra are featureless in the
L/M bands, so the observed spectral features are purely telluric. Each
standard is observed with an AB nod pair along the slit.

We reduce each nod pair with \esorex{} \texttt{cr2res\_obs\_nodding} and
then fit the extracted A and B spectra independently with \vipere{},
a fork of the \texttt{viper}\footnote{\url{https://github.com/mzechmeister/viper}}
tool \citep{viper} adapted for
CRIRES+\footnote{\url{https://git.astro.lavail.net/alexis/vipere}},
a forward-modelling tool that simultaneously fits the
telluric absorption, wavelength solution, continuum, and instrumental
profile. The key output is a wavelength polynomial per detector order and
per nod position.

The slit tilt at each pixel is then computed as the wavelength difference
between the A and B fits, converted to a pixel shift and divided by the
measured nod throw (in pixels). The nod throw is determined by fitting
double Gaussians to the spatial profile of the combined A frame.

Collecting measurements across all settings and chips, we fit a linear
relation between tilt and wavelength for each band:
\begin{align}
  t_L(\lambda) &= -6.97 \times 10^{-5}\,\lambda + 0.203 \quad
    (\mathrm{rms} = 0.005), \label{eq:tiltL} \\
  t_M(\lambda) &= -4.68 \times 10^{-5}\,\lambda + 0.169 \quad
    (\mathrm{rms} = 0.006), \label{eq:tiltM}
\end{align}
where $\lambda$ is in nm and $t$ is the tilt in pixels per spatial pixel.
The fit uses iterative 3$\sigma$ clipping and interpolates smoothly
across the CO$_2$ fundamental absorption band near 4.2--4.5\,\um{} where
tilt measurements are unreliable. The resulting tilt values are written
into the \texttt{SlitPolyB} column of the tracing tables used for all
subsequent reductions.

\begin{table}
\caption{Telluric standard stars used for slit tilt and wavelength
  calibration.}
\label{tab:tellstds}
\centering
\begin{tabular}{llll}
\hline\hline
Settings & Star & Sp.\ type & Slit \\
\hline
L3244--L3302, L3412, L3426 & $\alpha$\,Eri & B6\,V & 0\farcs2 \\
L3340 & HD\,153426 & O9\,II & 0\farcs4 \\
L3377 & $\theta$\,Aql & B9.5\,III & 0\farcs4 \\
M4187, M4266 & HR\,7950 & B & 0\farcs2 \\
M4211 & $o$\,Hya & B9\,III & 0\farcs2 \\
M4318 & $\theta$\,Aql & B9.5\,III & 0\farcs4 \\
M4368 & $\beta$\,Ori & B8\,Ia & 0\farcs2 \\
M4416--M4504 & $\alpha$\,Eri & B6\,V & 0\farcs2 \\
M4519 & V*\,S\,CrA & T\,Tauri & 0\farcs2 \\
\hline
\end{tabular}
\tablefoot{All standards are hot stars with featureless L/M-band continua,
  except M4519 where V*\,S\,CrA was the only available target; its
  \vipere{} fits are adequate but noisier.}
\end{table}


%--------------------------------------------------------------------
\subsection{Wavelength calibration}
\label{sec:wavecal}

The \vipere{} fits to the telluric standards also provide improved
wavelength solutions. For traces where the fit converged well
(percentage RMS residual $< 10$ and wavelength shift from the pipeline
solution $< 10$\,\kms), the \vipere{} polynomial replaces the pipeline
wavelength calibration in the tracing table. This applies to 225 out of
311 traces (72\%). The remaining traces, typically at the longest M-band
wavelengths where telluric features are weak or in the CO$_2$ gap, retain
the pipeline wavelength solution.


%--------------------------------------------------------------------
\subsection{Science extraction}
\label{sec:extraction}

\crires{} exhibits mechanical flexure that shifts the spectral orders on
the detector by typically 5--10 pixels between the daytime calibration
and the nighttime science observation, occasionally reaching up to
${\sim}50$ pixels. We correct for this by cross-correlating the spatial
edge profile of each raw science frame against the order boundaries
defined in the reference tracing table, measuring a global Y-shift that
is applied to the trace polynomials before extraction.

The adjusted tracing table, together with the nearest-in-time master flat
and blaze function from 233 calibration epochs, forms the input to
\esorex{} \texttt{cr2res\_obs\_nodding}, which performs optimal extraction
of the A and B nod spectra.


%--------------------------------------------------------------------
\subsection{Telluric correction}
\label{sec:tellcorr}

Each extracted science spectrum is fitted with \vipere{}, which models
the observed spectrum as a continuum polynomial multiplied by the
convolved atmospheric transmission, accounting for H$_2$O, CO$_2$, CH$_4$,
CO, O$_2$, and N$_2$O absorption. The telluric transmission model is
reconstructed from the fit residuals as
$T = (\mathrm{observed} - \mathrm{residual}) / C$, where $C$ is the
fitted continuum. The corrected spectrum is then
$S_\mathrm{corr} = S_\mathrm{obs} / T$.

For orders where the \vipere{} fit fails (e.g.\ the CO$_2$-saturated
region near 4.26\,\um), the original uncorrected spectrum is retained
with the telluric transmission column set to NaN.

Wavelengths in the output spectra are updated using the \vipere{} fit
where available. For unfitted orders, wavelengths are interpolated using a
2D polynomial surface fit (degree 2 in both pixel position and spectral
order number) to the \vipere{} wavelengths of the fitted orders on the
same chip.


%--------------------------------------------------------------------
\section{The archive}
\label{sec:archive}

The reduced spectra are served through a web application at
\url{https://neon.physics.uu.se/crires-lm/}. The front page presents a
filterable and sortable table of all observations for which telluric-corrected
spectra are available, currently 3913 nod pairs. Users can filter by target
name, wavelength setting, or ESO programme ID.

Each observation page shows an interactive spectrum plot (using Plotly)
displaying the telluric-corrected spectrum, the fitted telluric model, and
the continuum, alongside diagnostic plots of the extraction and telluric
fit quality. FITS files containing the corrected spectra are available for
download; the output columns per spectral order include the corrected
spectrum, error, wavelength, telluric transmission, and fitted continuum.

Individual AB pair reductions are accessible alongside the combined
template views. The full reduction code is available at
\url{https://github.com/ivh/CRIRES-LM}.


%--------------------------------------------------------------------
\section{Use of AI tools}
\label{sec:ai}

% TODO: write this section
Claude Opus 4.6 \citep{claude_opus46} was used ...


%--------------------------------------------------------------------
\section{Summary}
\label{sec:summary}

We have reprocessed all 11\,131 public \crires{} L- and M-band science
frames from the ESO archive into 5649 telluric-corrected nod-pair spectra
covering 156 targets across 16 wavelength settings. The reprocessing
includes slit tilt and wavelength calibration from telluric standard fits,
per-observation flexure correction, and forward-model telluric correction.
The resulting archive is browsable and downloadable at
\url{https://neon.physics.uu.se/crires-lm/}.

\begin{acknowledgements}
Based on data obtained from the ESO Science Archive Facility under
programmes listed at the archive URL. This work made use of the
\crtwores{} pipeline \citep{cr2res} and the \vipere{} telluric fitting
tool, a fork of \texttt{viper} \citep{viper}.
\end{acknowledgements}


%--------------------------------------------------------------------
\bibliographystyle{aa}
\bibliography{crires_lm}

\end{document}
